% Nouveau EGA V (cote 25): archive cote 25, site page 65 (PDF page 66).
% typescript. Draft compared: none.
% AI-transcribed by Claude Opus 5.5 (claude-opus-5-5, adaptive thinking, effort high), Message Batches API, batch ega25b (msgbatch_01A7n5YuXubUPNK1zsepi8RG), 2026-11-13; not yet checked by eye.
\scanpage{25/65}
% NOTE: the top right corner is cut off; a typed number, possibly "33", is only partly visible there
\mspage{\unsure{33}} \mspage{\hand{60}}

\unsure{permettant|permettamt} d'identifier \over{L}{\unsure{l}}es points de $\hand{\check{}}P$\hand{,} à valeurs dans \xout{xxxSi} S disons\hand{,}
comme les sous-variétés linéaires de codimension 1 de P (appelés enco-%
re \uline{hyperplans} de P). Ceci dit, $L^{o}$ consiste en les hyperplans qui
\uline{contiennent} la sous-\xout{\unsure{va}}riété linéaire L de P (par quoi bien s\over{u}{û}r on
entend que les points de $L^{o}$ \xout{xxxxx} à valeurs dans un S' sont les
hyperplans dans $P_{S'}$ qui contiennent $L_{S'}$ ). Ceci résulte du fait \xout{qxx}
\margin{|}(qui aurait dû venir en même temps que le fait qu'une \inserted[typed]{sous-}variété linéaire \inserted[typed]{L de P}\hand{,}
détermine le quotient loc libre $\underline{\unsure{G}}$ de $\underline{E}$ dont il provient) que si
$\underline{G}$ et $\underline{G}'$ sont deux quotients loc libres de $\underline{E}$ (pas néc de même rang)
alors $P(\underline{G}')\hand{\hookrightarrow}P(\underline{G})$ (en tant que sous-variétés linéaires de $P(\underline{E})$) ssi
$\underline{G}'$ est majoré par $\underline{G}$ (et l'inclusion $P(\underline{G}')\hand{\rightarrow}P(\underline{G})$ n'est alors autre
que le morphisme déduit de $\underline{G}\hand{\rightarrow}\underline{G}'$ ).

Voici un minimum de sorites dont on devra disposer. La liste
complète ne pourra de toutes façons être fixée qu'une fois que l'ensem-%
bles des autres n$^{\text{o}}$s du présent paragraphe sera\xout{\unsure{nt}} rédigé\xout{\unsure{s}}.

\deleted{12. \uline{Ratrappage de résultats \xout{xxbxix} sur les sections linéaires} .}

\hand{Il me semble commode d'introduire}
\hand{aussi le foncteur}
\[ \hand{\uline{Grass}_{\unsure{\bullet}}(\underline{E})(S') = \text{ens des Modules quotients}} \]
\hand{loc. libres (de rang}
\hand{non précisé) de $\underline{E}_{S'}$}

\hand{Alors \uline{Grass}$_{\unsure{*}}(\underline{E})$ est représentable par}
\hand{$\coprod_{n\geq 0}$ \uline{Grass}$_{n}(\underline{E})$. Les \uline{sous-variétés linéaires}}
\hand{de $P(\underline{E})$ sont en fait définis par les}
\hand{sections de \uline{Grass}$_{\unsure{*}}(\underline{E})$ sur S [NB. le}
\hand{rang i.e. la dim relative peut varier si S non connexe)}

\mspage{\hand{65}}
